% THETA3-B16 contribution by /root/definite_bridge, 21 November 2026.
\subsection{THETA3-B16: a pole criterion that also supplies finite coherence}\label{B16}

The value-at-one criteria in B15 start with an automorphic representation on a chosen smaller group. There is also a useful converse route: a pole on the rank-three group can construct a coherent smaller group. Its character conventions must be retained.

\paragraph{A precise conditional descent statement.}
Use exactly the notation of Jiang--Wu, \emph{On $(\chi,b)$-factors of cuspidal automorphic representations of unitary groups I}, arXiv:1409.0767v3, \S3.1 and Theorem 3.3(1), pp.15--17. Let $X$ be a skew-Hermitian space of rank three, let $\chi$ be a unitary Hecke character of $E$ with $\chi|_{\A_F^\times}=1$, and let $\sigma$ be irreducible and cuspidal on $\U(X)(\A_F)$. Write
\[
 \rho=\sigma\otimes\chi^{-1}
 \tag{B16a}
\]
in the source's notation for its twist on the unitary group. Suppose that:
\begin{enumerate}
\item $L^S(s,\sigma\times\chi)$, with the source's convention, has a pole at $s=1$;
\item $\rho$ has a real component which is not one dimensional;
\item at each real place, the only nonzero rank-two local theta lift of the corresponding input, with this splitting and additive character, is on a definite signature.
\end{enumerate}
Then there is a global Hermitian space $S$ of rank two, definite at every real place, for which the global theta lift of $\rho$ to $\U(S)$ is nonzero and cuspidal. In particular its finite local isometry classes are coherent, and its local theta lifts are nonzero at all finite places. No finite local form or lattice has been prescribed in this statement.

\textbf{Proof.}
In Jiang--Wu's Theorem 3.3(1), take their $m=3$, $\epsilon_\chi=0$, and $j=1$. Its distinguished point is
\[
 \tfrac12(m+1-\epsilon_\chi)-j=1,
\]
and it gives $\operatorname{LO}_{\psi,\chi}(\rho)\leq2$. The lowest occurrence minimizes the first-occurrence dimension over global Hermitian Witt towers of the prescribed even parity; see \S3.2, p.16. Therefore the dimension realizing the minimum is zero or two. At dimension zero the theta series is the determinant character $\nu_\chi\circ\det$, as the source states immediately before \S3.2. A nonzero global pairing with it would force every local component of the irreducible input to be one dimensional. Hypothesis 2 excludes this. Thus the lowest occurrence has dimension two, and by its definition its global theta lift is nonzero and cuspidal.

Nonzero global theta occurrence implies the corresponding local occurrences. Hypothesis 3 forces the real signatures of the constructed space to be definite. The theorem constructed a global space, so its finite local forms already satisfy the global coherence relation. This proves the assertion. \hfill$\square$

This result replaces separate finite-place guesses by one global pole test, once its representation and character dictionary has been established. It does not assert that the pole is known for the 032 wedge. In particular, the twist in (B16a) must be resolved against the bilinear integral in Jiang--Wu's \S3.1 and the contragredient in the local Howe quotient before substituting the canonical representation in (B15d). The local computations of B15 give the required unique definite occurrence after that convention matching; they do not by themselves prove hypothesis 1.

\paragraph{The invariant under a compensating splitting change.}
This remaining test is not removed by the independence of the special-unitary action from splitting characters. Let $\tau$ be a global character of $\operatorname{Res}^{1}_{E/F}\mathbf G_m$ and put
\[
 \tau_E(z)=\tau(z/\overline z),\qquad
 \Pi'=\Pi\otimes(\tau\circ\det),\qquad
 \chi'=\chi\tau_E.
\]
In the standard local-parameter convention of (B15c), changing the splitting by $\tau_E$ multiplies the input action by $\tau\circ\det$. Hence the compensating simultaneous twist transports the same local Howe correspondence on the same underlying Hermitian spaces. This is the global version of (B15g). Jiang--Wu, equation (3.1), p.15, gives the same determinant-character rule, with $\nu_\kappa(z/\bar z)=\kappa(z)$.

There is an equally explicit invariant of unramified standard parameters. Whenever standard base change is used to write them, one has
\[
 \operatorname{BC}(\Pi')\otimes(\chi')^{-1}
   =\operatorname{BC}(\Pi)\otimes\chi^{-1}.
 \tag{B16b}
\]
Indeed, twisting a representation by $\tau\circ\det$ tensors its standard base-change parameter by $\tau_E$; the two additional factors then cancel. This also proves equality of the corresponding incomplete standard $L$-functions term by term at every unramified place. Formula (B16b) is stated in this parameter convention; it is not an unexplained replacement of Jiang--Wu's named $L$-function by a differently twisted one.

To change the relative splitting further by a character $\eta$ changes the right side of (B16b) by $\eta^{-1}$. That is a different standard parameter and a different pole question. The original nonzero A7 lift supplies information for its original source $P$ and characters. The compensating transport preserves $P$, including any real signature $(1,1)$; an additional relative change has not inherited that nonvanishing. This identifies the missing representation-theoretic calculation without asserting that another definite realization is impossible.

Finally, a successful application to each of several rank-three representations separately need not produce the same finite local source forms. The \texttt{Comm} in B13 is common to the whole family. Its use still requires common spaces and real action across tuples, the specified finite lattice data, and the kernel/covolume equation (B15m). The pole construction therefore closes a single-representation coherence problem under its hypotheses; it does not supply the entire arithmetic family by itself.
